\section{My support}

\subsection*{4a. What others did that helped me work more effectively}

Edward is the reason the Cube spoke to the Pi at all. In wk15 I spent most of a weekend trying to get pymavlink to read from \texttt{/dev/ttyAMA0} on Python 3.13, and it was Edward who pointed out --- after I had filed three dead-end bug reports in our Teams channel --- that \texttt{pyserial} serial reads were broken on Python 3.13 and we should route through a \texttt{mavproxy} UDP bridge. That was a 40-minute conversation that saved me probably two more days of solo debugging. I recorded the exact command in \texttt{MEMORY.md} so the team would not rediscover the workaround. I am naming this specifically because in the moment I resented the interruption --- I had a plan, I was working, I did not want to be told my plan was wrong --- and in retrospect the interruption was the single most valuable technical contribution anyone made to my side of the codebase. The uncomfortable realisation was that my instinct to ``finish debugging before asking for help'' is a cost I pay, not a professionalism I exhibit.

Robin helped me by pushing back on the state machine in wk17. I have already described that disagreement in §2b, so I will not re-narrate it here, but I want to explicitly credit him: without his pushback I would have shipped a 32-state FSM that only I could maintain, and the integration of his flight script with my vision output (wk20) would have been a nightmare. The feedback he gave me was not comfortable to hear. It was also correct. I want to record that in writing because giving and receiving unflinching technical feedback is exactly the thing M16/M17 asks us to evidence, and the normal thing in student projects is for those exchanges to be smoothed over into silence.

Demetro helped me in a different way: by being willing to let me write his slides and speaking scripts without turning it into a status contest. When he asked for help with his two Final Design Review slides (\texttt{32fb0c5}) and then with the speaking scripts (\texttt{b2a9cd4}), he was explicit about what he needed and willing to accept drafts that I had written in my voice rather than his. That made the collaboration fast --- I could build what I knew would land with a client audience without negotiating every sentence. Not every teammate would have been that generous with authorship, and it is worth recording that the generosity was his, not mine.

Sid --- the technical assistant, not a teammate --- deserves a specific mention. He provided the single DJI flight video I used for FOV calibration and dataset generation. It was one video. It arrived late. I used it for everything: \texttt{tools/fov\_calibrate\_video.py} (clicking the dummy at 9 altitudes), \texttt{generate\_dataset\_v2.py} (background frames for 300 synthetic composites), \texttt{tools/label\_tool.py} (16 hand-labelled real frames), and the A/B model comparison in \texttt{video\_test\_compare.py}. The retrained model (mAP50 = 0.995) does not exist without that video. I want to name the dependency because it is a case study in how a single external data point from outside the team can unblock a significant chunk of work, and the correct response is to use it ruthlessly and acknowledge it honestly.

\subsection*{4b. What I did that helped others work more effectively}

I am going to be specific rather than comprehensive, because the exemplar guide is explicit that generic claims of ``I helped the team'' are automatic Merit-ceilings.

I built Robin's integration path. When it became clear in wk19 that Robin's flight script needed detection results from my vision code without inheriting my whole codebase, I spent the wk20 weekend writing \texttt{passive\_watch\_2.py} with a dedicated \texttt{--robin} flag (\texttt{a5b1ab7}), an HTTP \texttt{cv\_mode} polling endpoint, and a two-image output format matched to what his state machine expected. I then wrote him a parsing-example README (\texttt{885cee4}) with the exact CLI invocation and the JSON schema laid out in copy-pasteable form. The effect was concrete: Robin got his integration working in one afternoon instead of the three days he had estimated. The mechanism that made this work was anticipating his API needs before he asked --- I read his state machine on a Sunday evening and wrote the integration around his calling convention, not mine. That inversion (caller-shaped rather than author-shaped) is a habit I want to keep.

I wrote Demetro's slides and speaking scripts. When he asked for help with his two Final Design Review slides (\texttt{32fb0c5}), I built them from scratch rather than sending him bullet points --- full-screen layouts, inline diagrams, and then iterated 4--5 times on the speaking scripts (\texttt{b2a9cd4}, \texttt{ba667f2}, \texttt{9b826f8}, \texttt{7920e5a}, \texttt{794ef5c}) until the delivery timed out at 3.1 minutes with an inline teleprompter sidebar. The speaking scripts were the piece he had not asked for and the piece he needed most, because he told me afterwards he had been nervous about the wording. I do not count that as leadership --- it is a favour --- but I count it as evidence that I was paying attention to what a teammate was scared of rather than what he was asking for.

I drafted the complaint letter on behalf of the team (\texttt{docs/COMPLAINT\_LETTER.md}). I have already described why I think this was the hardest and most important thing I did in the project (§1b). What I want to add here, in the ``helping others'' frame, is the specific observation that my first three drafts were written in my voice and my anger, and the fourth was rewritten in the team's voice after I circulated it for comments. The pattern I noticed in myself is that when I write under emotional load I default to first-person-singular prose, and for a team letter the correct register is first-person-plural and past-tense. Catching that in myself required showing the draft to someone else. I will not forget the specific correction: Robin suggested cutting three sentences that were ``technically accurate but personally sharp,'' and he was right. The final letter is measured because two other people edited my first-person anger out of it. That is the purest example I have from this project of feedback improving a deliverable I thought was already done, and it is the main reason I will not submit any important external-facing writing as a solo draft again.

I also want to note one specific piece of \textbf{feedback I gave to Robin} around wk20. His flight script had no timeout guards on the \texttt{CENTERING} state --- if the target was lost mid-descent, the loop would hang waiting for a detection that never came. I told him this in a Teams message, named the specific failure mode (``consecutive\_lost counter doesn't reset''), and suggested the fix (time-based timeout with fallback to ascend-and-retry). His reaction was to fix it the same afternoon (the fix is now in \texttt{4\_detect\_and\_center.py}), and a week later in wk21 he pulled me aside after the team meeting to say he had caught a similar timeout bug in a separate module because he now ``looked for that shape of bug.'' That is the outcome I want from feedback --- not that the specific bug was fixed, but that the shape of the check propagated into his working model. I delivered the message well because I named the exact variable and the exact line rather than the general concern, and I now see that the specificity is what made it feel like a collaboration instead of a criticism. The generalisable lesson I am taking forward is: when giving feedback to a high-performing peer, the scarce resource is specificity, not softness.

\subsection*{Closing}

I entered this module thinking that engineering leadership is technical excellence plus the discipline to organise your own work. I leave it thinking it is mostly the skill of making other people's work easier, and of making your own disagreements visible before they become expensive. The specific moment this became clear to me was the wk17 FSM argument with Robin --- the one in which I was partly right on the technical merits and fully wrong on the team fit. I am still uneasy about the leftover work I did not finish: I never paired with Edward on the hardware side, the \texttt{main.py} state machine remains an unmaintainable 32-state graph that only I can read, and the 12/12 requirements-in-simulation / 4/12 requirements-on-hardware split is a gap my code cannot close. Those are real failures of impact that I cannot claim credit for having solved. The lesson I carry is that I can convert blocked dependencies into buildable work almost reflexively --- that is my strength --- and that the same reflex is what makes me a solo-worker in a team where solo work is the expensive failure mode. Next project I will use the first instinct and watch for the second.
